% PDF 30
\sourcepage{26}{30}
函子 $\mathbb E^I$ 由 $k$-代数 $P_I$ 表示\footnote{初学者从这里所需的具体情形就能看出可表函子的概念何等简单：$\mathbb E^I\colon\G_k\to\Sets$，且 $k'\mapsto\mathbb E^I(k')=\Hom_{\G_k}(P_I,k')$；每个 $k$-代数态射 $u\colon k'\to k''$ 都以显然的协变方式诱导 $\mathbb E^I(u)\colon\Hom(P_I,k')\to\Hom(P_I,k'')$。显然，我们并未使用范畴 $\G_k$ 的任何特殊性质。可表函子的主要定义见第 1 节。}。
对于每个 $k'\in\Ob\G_k$，$\mathbb E^I(k')$ 中的点称为 $\mathbb E^I$ 的 $k'$ 值点。

现在设有一族多项式方程
\begin{equation}\tag{5.3} f_j(x)=0,\qquad j\in J,\end{equation}
其中 $J$ 是任意指标集，$f_j(T_i)\in P_I$（见第 1 节），$j\in J$。让 $k'\in\Ob\G_k$ 变化，考察 $\mathbb E^I(k')$ 中满足 $S$ 的所有点组成的集合 $V_S(k')\hookrightarrow\mathbb E^I(k')$，就得到函子 $k'\mapsto\mathbb E^I(k')=k'^I$ 的子函子 $k'\mapsto V_S(k')$。

\begin{quote}\bfseries 从代数几何的角度研究方程族 $S$，就是研究函子 $V_S$。\end{quote}

即使在经典情形下，$k$ 是域，例如实数域 $\mathbb R$，而方程及不定元的个数都是有限的，集合 $V_S(\mathbb R)$ 仍可能为空。例如一元方程 $x^2+1=0$，或实平面中的方程 $x^2+y^2=0$\footnote{译注：此处原文如此。该方程有实解 $(0,0)$；结合下文“虚圆”，作者所指应是 $x^2+y^2+1=0$。}。然而，我们并不因此认为这些方程所定义的“簇”是平凡的，因为在适当的 $\mathbb R$-代数上，这些方程是有解的。例如取 $k'=\mathbb C\in\G_{\mathbb R}$，即复数域，第一种情形得到两个解 $\pm i$，第二种情形则得到所谓虚圆上的一整个连续统的复值点。因此，如果只局限于实数，而不允许 $k'$ 变化，关于方程组 $S$ 实际上就几乎得不到什么信息。
% PDF 31
\sourcepage{27}{31}
这里 $k'\in\Ob\G_k$。\textbf{几何性质应当在任意基变换下保持不变}；前面例子中“$V_S(\mathbb R)$ 为空”这一性质就不具有这种不变性。

本节末尾将补充若干历史背景，说明为什么要引入 $\mathbb E^I$、$V_S$，而不把自己限制在 $k$，特别是 $\mathbb R$、$\mathbb C$ 等之上。

首先，定义 \textbf{$k$ 上的 $I$ 型标准仿射线性空间}\footnote{$\mathbb E^I$ 很快将作为一般仿射代数空间的一个特例出现；因此这里加上“线性”二字，尽管“仿射线性”乍看似有重复。这一用语与通常区分仿射代数簇和仿射空间的习惯一致。Grothendieck 常称之为 $I$ 型标准仿射空间。}为一个从 $k$-代数范畴 $\G_k$ 到集合范畴的函子 $\mathbb E^I$：它把每个 $k'\in\Ob\G_k$，即每个 $k$-代数 $k'$，对应到仿射空间 $k'^I$。也就是说，
\begin{align}\tag{5.4}\mathbb E^I&\colon\G_k\longrightarrow\Sets,\\\tag{5.5}\mathbb E^I&\colon k'\longmapsto k'^I\qquad(k'\in\Ob\G_k).
\end{align}
现在考虑 $\G_k$ 中的态射，即 $k$-代数同态（见第 2 节）
\begin{equation}\tag{5.6}u\colon k'\longrightarrow k''.
\end{equation}
它诱导映射
\begin{equation}\tag{5.7}\mathbb E^I(u)\colon k'^I\longrightarrow k''^I,\qquad(x_i)_{i\in I}\longmapsto(u(x_i))_{i\in I}.
\end{equation}
所以 $\mathbb E^I$ 的确是函子 $\G_k\to\Sets$。

\textbf{函子 $V_S$。}现在将系数属于 $k$ 的方程组 (5.1) $S$ 看作系数属于 $k'$ 的方程组，考察它的解集。
% PDF 32
\sourcepage{28}{32}
\begin{equation}\tag{5.8}V_S(k')=\{x\in k'^I\mid f_j(x)=0,\ \forall j\in J\}.
\end{equation}
于是，对于从 $k'$ 到 $k''$ 的每个同态 (3.4)，$k'^I$ 中满足 $S$ 的任一点，其在 $k''^I$ 中的像都满足同一组方程，即 $V_S(k')$ 被映入 $V_S(k'')$。用函子的语言说，范畴 $\G_k$ 中的 $k'\xrightarrow{u}k''$ 对应于集合范畴中的
\begin{equation}\tag{5.9}V_S(u)\colon V_S(k')\longrightarrow V_S(k'').
\end{equation}
换言之，\textbf{方程组 $S$ 定义了从 $\G_k$ 到集合范畴的函子 $\mathbb E^I$ 的一个子函子 $V_S$}：
\begin{equation}\tag{5.10} k'\xmapsto{V_S}\{x\in k'^I\mid f_j(x)=0,\ \forall j\in J\}\subset\mathbb E^I(k')=k'^I.
\end{equation}
即 $V_S(k')$ 是 $S$ 的 $k'$ 值解的集合。集合 $V_S(k')$ 称为仿射空间 $\mathbb E^I(k')$ 的一个\textbf{闭代数子集}\footnote{稍后将从仿射空间 $k'^I$ 上所谓 Zariski 拓扑的角度说明这一术语的来由，见本章第 14 节。}。

\textbf{［例子与历史动机。}算术问题，即求出 $V_S(k)$，本身可能极为困难；除此以外，我们还看到 $V_S(k)$ 有时会为空，例如 $\mathbb R^n$ 中的虚二次曲面 $\sum_{i=1}^n x_i^2+1=0$。包含映射 $\mathbb R\hookrightarrow\mathbb C$ 对应于 $\varnothing\hookrightarrow V_S(\mathbb C)$。显然，从实数域过渡到复数域的做法由来已久，而且十分常见。许多代数和算术问题都要求扩张基域 $k$（参见 Dieudonné，\textit{Advances}，I）。于是出现了一个新问题：如何比较坐标属于不同扩域 $k'$、$k''$ 的点？如果二者不是同一第三个域的子域，就根本无从比较。A.~Weil 引入了泛域 $\Omega$，即在其素域上具有无限超越次数的域，并假设所需的一切扩域都是 $\Omega$ 的子域。Grothendieck 的方法不要求基环为域，也不采用 $\Omega$。相应地，他让 $k'$ 变化，考察不同的 $V_S(k')$；采纳 Kähler 的建议，这些 $k'$ 是 $k$-代数。他由此系统地研究函子 $V_S$，不再只研究固定 $k$ 时由 $S$ 表示的单一集合 $V_S(k)$，即所谓“代数簇”。

在几何上，$k$-代数起着 $V_S$ 上“$k$ 值函数”的作用。（$k$-代数同态所给的点与坐标在 $k$ 中的“常值”点一样，都满足 $S$。）经典代数几何早已使用这些代数，并称之为“坐标环”：$k'=k[\xi_1,\ldots,\xi_n]$。
% PDF 33
\sourcepage{29}{33}
在这里，考察的是仿射不可约代数簇，基环 $k$ 特别取为域，例如 $k=\mathbb C$，且 $I=\{1,2,\ldots,n\}$ 为有限集。这时 $k[\xi_1,\ldots,\xi_n]$ 是有限生成整环；$\xi_1,\xi_2,\ldots,\xi_n$ 是 van der Waerden 所说的“一般点”（\textit{allgemeiner Punkt}）的坐标。$k'$ 同构于多项式环 $k[x_1,\ldots,x_n]$ 对某个素理想 $\mathfrak J$ 的商环。

我们将看到，在 Grothendieck 的一般情形下，每个 $k'$ 都可以写成 $P_I/\mathfrak J$，其中 $P_I$ 是任意一族不定元 $(T_i)_{i\in I}$ 的多项式环，而 $\mathfrak J$ 是任意理想，不必为素理想。固定 $I$ 时，$P_I$ 的不同理想对应不同的函子。如果只考察 $k$ 值点，这一说法便不成立。例如，对每个 $m\geq1$，方程 $x^2+y^2-1=0$、$(x^2+y^2-1)^2=0$ 和 $(x^2+y^2-1)^m=0$ 都有相同的复值解集，所以主理想 $(x^2+y^2-1)^m$（$m=1,2,\ldots$）全都定义同一个 $V_S(\mathbb C)$。参见下一节。\textbf{］}

\Section{6}{对 $V_S$ 的内蕴研究与嵌入 $V_S\hookrightarrow\mathbb E^I$}
\textbf{提要。}我们的代数几何问题已归结为研究 $\mathbb E^I$ 的子函子 $V_S$。这里要区分两个问题：a）$V_S$ 的内蕴研究；b）嵌入 $V_S\hookrightarrow\mathbb E^I$ 的研究。这一区分，与初等微分几何中区分曲线或曲面的内蕴性质和依赖其嵌入的性质相对应：前者与嵌入三维空间或 $n$ 维空间（$n\geq3$）的方式无关，后者则同时依赖于曲线、曲面本身及其嵌入。

对子函子 $V_S$ 可以采取两种不同的观点：a）首先，内蕴地研究函子 $V_S$，即在相差同构的意义下考察它；这只是研究一个从 $k$-代数范畴 $\G_k$ 到集合范畴的函子。因此可以比较 $V_S$ 与 $V_{S'}$，其中 $S'$ 是在同一基环 $k$ 上定义的另一组方程，它的指标集 $I'$、$J'$ 可以不同。b）其次，研究嵌入 $V_S\hookrightarrow\mathbb E^I$ 的特殊性质。

\textbf{例。}\footnote{在代数平面曲线的非内蕴研究中，极曲线、切线、法线等都是确实需要使用曲线外部的点的经典例子。内蕴研究则更进一步，也更为精细。}观点 a）的一个例子，是研究处于某个有限维仿射空间中的实或复代数簇，而不依赖它们的嵌入方式；这样，即使环境维数 $N\ne N'$，也可以通过某种等价关系把特定的一对仿射簇视为相同。读者不妨想到嵌入实或复三维空间的代数曲线或曲面。当我们把这些簇的某个具体嵌入也考虑在内时，便得到了观点 b）的例子。
% PDF 34
\sourcepage{30}{34}
新意在于函子的观点：$k'$ 是变化的！这与初等微分几何或经典代数几何中固定 $\mathbb R$ 或 $\mathbb C$ 的做法不同。至于内蕴性质和外在性质的区分，则早已有之。

\Section{7}{化归到 $S$ 是 $P_I$ 的理想的情形}
\textbf{提要。［}两个系数在 $k$ 中、分别具有指标类型 $(I,J)$、$(I',J')$ 的多项式方程组 $S$、$S'$，称为等价的，当且仅当它们的解函子 $V_S,V_{S'}\colon\G_k\to\Sets$ 等价\footnote{即对每个 $k'\in\Ob\G_k$，存在集合之间的等价 $h(k')\colon V_S(k')\xrightarrow{\sim}V_{S'}(k')$。}。

固定 $I$，也就固定了“环境空间”$\mathbb E^I$。我们可以给 $S$ 添加所有有限线性组合 $\sum_{j\in J}g_jf_j$，得到 $P_I$ 中的理想 $\mathfrak J=P_I\cdot S$。于是 $V_S\sim V_{\mathfrak J}$，所以只需研究形如 $V_{\mathfrak J}$ 的函子，其中 $\mathfrak J$ 是 $P_I$ 的理想。此外，若 $\mathfrak J\ne\mathfrak J'$，则函子 $V_{\mathfrak J}$、$V_{\mathfrak J'}$ 不等价\footnote{这说明，即使 $I$ 为无限集而 $k$ 为交换域，也不宜把 $S$ 限制为有限集：对 $S$ 最自然的假设，是它为 $P_I$ 的理想，而一个理想作为集合几乎从不有限。不过在经典代数几何中，$k$ 为域且 $I$ 有限，由 $P_I$ 中理想的 Hilbert 基定理，总可以用有限的 $S$ 来定义同一个对象。即便如此，要求 $I$ 有限，乍看仍是人为且会带来技术上的不便。}\footnote{译注：这里以及下文的一一对应，须在 $\mathbb E^I$ 的子函子意义下理解，即保留给定嵌入；若仅比较抽象函子的同构类型，不同理想也可能给出同构的函子。}。换言之，我们将证明：能够由多项式方程组描述的 $\mathbb E^I$ 的子函子，与 $P_I$ 的所有理想之间有一一对应。\textbf{］}

无论从算术还是几何出发，首先都会注意到，可以用许多方式扩充方程组 $S$，而不改变它在基环 $k$ 或任意 $k$-代数 $k'$ 中的解。具体地说，任取线性组合 $\sum_{j\in J}g_jf_j$，其中 $g_j\in P_I$，只有有限多个 $g_j$ 非零，则 $S$ 的任一解 $x$ 也是 $\sum g_jf_j=0$ 的解。

因此，可以把所有有限线性组合 $\sum_{j\in J}g_jf_j$ 都加入 $S$。这些线性组合恰好组成由 $S$ 生成的理想 $\mathfrak J=\mathfrak J_S=P_I\cdot S$。于是对每个 $k$-代数 $k'$，都有
% PDF 35
\sourcepage{31}{35}
\begin{equation}\tag{7.1}\{x\in k'^I\mid f_j(x)=0,\ \forall j\in J\}=\{x\in k'^I\mid f(x)=0,\ \forall f\in\mathfrak J\}.
\end{equation}
事实上，由 $S\subset\mathfrak J$，得到 $V_{\mathfrak J}\subset V_S$；反过来，前面的说明又给出 $V_S\subset V_{\mathfrak J}$。所以两个函子 $V_S$、$V_{\mathfrak J}\colon\G_k\to\Sets$ 等价：
\begin{equation}\tag{7.2}V_{\mathfrak J}\sim V_S.
\end{equation}
许多不同的方程集合 $S$，可以对每个 $k'\in\Ob\G_k$ 都定义相同的解集，因此选择环 $P_I$ 中的理想 $\mathfrak J$ 来定义解集 $V_{\mathfrak J}(k')$，是最自然的做法。我们马上就会看到，这还是一个典范的选择。［参见原书第 30 页脚注 (2)。］

下面进一步说明标准仿射空间 $\mathbb E^I$ 的闭代数子空间与 $P_I$ 的理想之间的对应。对每个理想 $\mathfrak J\subset P_I$，赋予子函子 $V_{\mathfrak J}\hookrightarrow\mathbb E^I$\footnote{一个 $k'$ 值点 $u\in\Hom_{k\text{-alg}}(A,k')$，被映到子集 $V_{\mathfrak J}(k')\subset\mathbb E^I(k')$ 的一个点：$u\mapsto(u(a_i))_{i\in I}$。$V_{\mathfrak J}(k')$ 是 $\mathbb E^I(k')$ 中满足方程组 $f(x)=0$（所有 $f\in\mathfrak J$）的全部点。}。我们断言，仅由 $V_{\mathfrak J}$ 就能重建理想 $\mathfrak J$。实际上，存在一一对应
\begin{equation}\tag{7.3}\mathfrak J\longleftrightarrow V_{\mathfrak J}
\end{equation}
把理想 $\mathfrak J\subset P_I$ 与上述类型的 $\mathbb E^I$ 的子函子联系起来。那么，怎样由 $V_{\mathfrak J}$ 恢复 $\mathfrak J$ 呢？

元素 $f\in P_I$ 属于 $\mathfrak J$，当且仅当对每个 $k'\in\Ob\G_k$ 以及 $V_{\mathfrak J}$ 的每个 $k'$ 值点 $x$，即 $x\in V_{\mathfrak J}(k')$（参见第 1 节），都有
% PDF 36
\sourcepage{32}{36}
\begin{equation}\tag{7.4}f(x)=0.
\end{equation}
事实上，如果 $f\in\mathfrak J$，显然满足 (7.4)。反过来也成立：若 $f\in P_I$ 对任意 $k'\in\Ob\G_k$ 以及 $V_{\mathfrak J}(k')$ 的任意 $k'$ 值点都满足 (7.4)，那么 $f\in\mathfrak J$。因为可以取 $k'=P_I/\mathfrak J$；对由典范同态 $P_I\to A=P_I/\mathfrak J$ 给出的点 $x$，$f(x)=0$ 当且仅当 $f\in\mathfrak J$，因为 $\mathfrak J$ 正是这一同态的核。

此外，\textbf{映射 $\mathfrak J\mapsto V_{\mathfrak J}$ 反转包含次序}。设 $\mathfrak J\subset\mathfrak J'\subset P_I$，那么有反向的包含关系 $V_{\mathfrak J'}\subset V_{\mathfrak J}$。反过来，根据刚才对 $\mathfrak J$ 的刻画，$V_{\mathfrak J'}\subset V_{\mathfrak J}$ 也蕴含 $\mathfrak J\subset\mathfrak J'$。换言之，映射 $\mathfrak J\mapsto V_{\mathfrak J}$ 为双射，因为它在两个方向上都保持次序\footnote{译注：原文此处写作“保持次序”；结合前文，准确地说应为两个方向上均反转包含次序。}。

于是，$P_I$ 中理想的语言，与子函子 $V_{\mathfrak J}\hookrightarrow\mathbb E^I$，即“闭代数子空间”的语言之间，得到了一部完整的词典。因此，如果坚持研究嵌入 $\mathbb E^I$ 的仿射代数空间，这项研究就等价于研究 $P_I$ 的理想。这也就是为什么“老一辈”会说，仿射代数几何与 $P_I$ 中的理想论“完全一样”。［第 8 节将讨论这种非内蕴观点的局限。］

\textbf{历史评注。}我们已经指出（见原书第 6 页的课程提要），经典代数几何中 $P_I=k[T_1,T_2,\ldots,T_n]$，其中 $n=\#I<\infty$。若 $k$ 代数闭，例如 $k=\mathbb C$，最值得关注的是根理想 $\mathfrak J$，即 $\mathfrak J=\sqrt{\mathfrak J}$（参见原书第 6 页脚注 (5)）。事实上，由 Hilbert 零点定理，$A=k[T_1,\ldots,T_n]/\mathfrak J$ 是某个代数簇 $M$ 上多项式函数 $f\colon M\to k$ 的环，当且仅当 $\mathfrak J=\sqrt{\mathfrak J}$。\footnote{原书此页再次刊有前页同一脚注：一个 $k'$ 值点 $u\in\Hom_{k\text{-alg}}(A,k')$ 被映到 $V_{\mathfrak J}(k')\subset\mathbb E^I(k')$ 中的点 $(u(a_i))_{i\in I}$；$V_{\mathfrak J}(k')$ 是满足全部方程 $f(x)=0$（$f\in\mathfrak J$）的点集。}
% PDF 37
\sourcepage{33}{37}
这里的 $M$ 不必不可约\footnote{$M$ 不可约，当且仅当 $\mathfrak J=\sqrt{\mathfrak J}$ 是素理想。}。反之，如果 $\sqrt{\mathfrak J}$ 严格大于 $\mathfrak J$，那么 $A$ 中就有非零幂零元 $\bar f\ne0$，即存在整数 $m>1$ 使得 $\bar f^{\,m}=0$。显然，$\bar f$ 不可能表示由 $\mathfrak J$ 表示的代数簇 $M$ 上的任何函数\footnote{任意 $\mathfrak J\subset k[T_1,\ldots,T_n]$ 都表示一个簇 $M_{\mathfrak J}$，但不同的理想可能表示同一 $M$。不过，在表示 $M_{\mathfrak J}$ 的理想中有一个最大的，恰好是 $\sqrt{\mathfrak J}$；参见所引书第 127 节，第 81 页。}。尽管面临这种“似乎不被允许”的情形，Gröbner 仍曾尝试\footnote{Gröbner，\textit{Moderne algebraische Geometrie (Die idealtheoretischen Grundlagen)}［《现代代数几何：理想论基础》］，Springer，维也纳与因斯布鲁克，1949。}在任意理想与 $k^n$ 的仿射代数簇之间建立双射 $\mathfrak J\leftrightarrow V_{\mathfrak J}$，却没有顾及 $\mathfrak J$ 和 $\sqrt{\mathfrak J}$ 定义相同的簇——这里的簇是指 $k$ 值点的处所，这也是当时唯一允许考虑的几何点。因此，他被批评为“同义反复”（例如 Samuel\footnote{Samuel，\textit{Méthodes d'algèbre abstraite en géométrie algébrique}［《代数几何中的抽象代数方法》］，Ergebnisse der Mathematik，第 4 卷，第 2 版，1967，第 128 页：“至于 W.~Gröbner……他的‘代数簇’看起来几乎就是多项式理想。”原文：\textit{Quand à W.~Gröbner\ldots, tandis que ses «Algebraische Mannigfaltigkeiten» ont l'air d'être à peu près des idéaux de polynômes.}}）。其实，只有引入谱 $X=\Spec A$ 的结构层 $\mathcal O_X$（参见第三章），这种同义反复才会完全消失。这时 $A$ 就成为 $\mathcal O_X$ 的整体截面环；而无论是否存在非平凡的幂零元，总能定义 $\mathcal O_X$。

\Section{8}{$k$ 上仿射代数空间的范畴；“绝对”情形 $k=\mathbb Z$}
\textbf{Grothendieck 的提要。}谈到老一辈的观点［参见第 6 节提要］，从我们正在采用的内蕴观点出发，真正感兴趣的并不是把代数空间嵌入某个特定的标准仿射空间。因此，理想 $\mathfrak J$ 本身并非首要对象；有意义的是商代数 $A=P_I/\mathfrak J$\footnote{请回顾第 4 节末尾所述，在经典情形下 $P_I/\mathfrak J$ 的几何意义。}，因为这个商将表示我们要研究的函子。准确地说，我们将看到：与某个 $V_S$ 同构的函子 $\G_k\to\Sets$，恰好是可表函子\footnote{参见第 1 节。在 $A=P_I$ 的特例中，我们已遇到可表函子：$P_I$ 表示标准仿射空间 $\mathbb E^I$。}$\mathfrak X_A$。这就是说，存在一个 $k$-代数 $A$，称为表示对象\footnote{$A$ 在相差同构的意义下唯一确定（参见第二章第 1 节）。}，使得 (8.2) 成立。

用 $P_I$ 的理想来表述，就是定义 $A=P_I/\mathfrak J$，其中 $\mathfrak J=P_I\cdot S$，于是有函子 $\G_k\to\Sets$ 的等价 $V_S\sim V_{\mathfrak J}\sim\mathfrak X_A$。$\mathfrak X_A$ 的优点在于，它不依赖于嵌入任何标准仿射空间的具体方式。反过来，每个 $k$-代数 $A$ 都定义一个 $\mathfrak X_A$，并与某个 $V_S$ 等价，因为总可以适当地选择 $I$ 与 $\mathfrak J$，使 $A$ 成为商 $P_I/\mathfrak J$。函子 $\mathfrak X_A$（$A\in\Ob\G_k$）称为\textbf{由 $A$ 表示的 $k$ 上仿射代数空间}。

$k$ 上的仿射代数空间组成一个范畴，它与 $k$-代数范畴 $\G_k$ 反等价\footnote{即与 $\G_k$ 的反范畴 $\G_k^{\op}$ 等价。}。

取 $k=\mathbb Z$，就得到“绝对”仿射代数空间范畴，它与含幺交换环的范畴 $\G=\G_{\mathbb Z}$ 反等价。
% PDF 38
\sourcepage{34}{38}
说函子 $V_{\mathfrak J}$（$\mathfrak J$ 为 $P_I$ 的理想）由 $A=P_I/\mathfrak J$ 表示，等价于叙述典范同态 $p\colon P_I\to A$ 的下列泛性质。设 $x\colon P_I\to k'$ 是 $P_I$ 的一个 $k'$ 值点，其中 $k'\in\Ob\G_k$（见第 2 节）。那么，对每个 $f\in\mathfrak J$ 都有 $f(x)=0$，当且仅当 $x$ 唯一地经由 $p$ 分解，即存在交换图
\begin{equation}\tag{8.1}
\begin{tikzcd}[ampersand replacement=\&]P_I\arrow[r,"x"]\arrow[d,"p"']\&k'\\A\arrow[ur,"x_A"']\&\end{tikzcd}
\end{equation}
也就是 $x=x_A\circ p$，其中 $x_A$ 由 $x$ 唯一确定。

映射 $x\mapsto x_A$ 给出 $\mathbb E^I$ 中在 $\mathfrak J$ 上消失的 $k'$ 值点的集合 $V_{\mathfrak J}(k')$ 与集合 $\Hom_{\G_k}(A,k')$ 之间的双射。证毕。换言之，\textbf{函子 $V_{\mathfrak J}\colon\G_k\to\Sets$ 由 $k$-代数 $A=P_I/\mathfrak J$ 表示}。

由
\begin{equation}\tag{8.2}k'\longmapsto\mathfrak X_A(k')=\Hom_{\G_k}(A,k'),\qquad k'\in\Ob\G_k
\end{equation}
定义的协变可表函子 $\mathfrak X_A\colon\G_k\to\Sets$，称为\textbf{由 $k$-代数 $A$ 表示的 $k$ 上仿射代数空间}。

前面的构造消去了 $\mathbb E^I$。表面上看，$A$ 似乎受到限制，因为它由像 $a_i=p(T_i)$（$i\in I$）生成。实际上没有这种限制：总能给 $A$ 选择一族 $k$-代数生成元，例如取 $A$ 的全部底集，然后把这一过程反过来做。用 Grothendieck 的原话说：我们将看到，给出这样的嵌入 $\mathfrak X_A\hookrightarrow\mathbb E^I$，作为从 $\mathfrak X_A$ 到 $\mathbb E^I$ 的态射，就必须考虑表示它们的代数之间方向相反的态射 $P_I\to A$。这又等价于给出 $P_I$ 的生成元 $T_i$（$i\in I$）的像 $a_i$。
% PDF 39
\sourcepage{35}{39}
这意味着，由 $T_i\mapsto a_i$（$i\in I$）唯一确定的同态 $P_I\to A$ 是满同态；而这又等价于 $A$ 是 $P_I$ 对某个理想 $\mathfrak J$ 的商。

这样就完全解决了如何不依赖任何仿射嵌入来刻画函子 $V_S\colon\G_k\to\Sets$ 的问题：

\begin{quote}\bfseries 函子 $V_S\colon\G_k\to\Sets$ 描述某个系数在 $k$ 中的多项式方程组 $S$ 的解，当且仅当 $V_S$ 可表，即存在 $A\in\Ob\G_k$ 使 $V_S$ 与 $\mathfrak X_A$ 等价。\end{quote}
反过来，对于每个 $k$-代数 $A\in\Ob\G_k$\footnote{在经典代数几何中，$k$ 是域，$A$ 是有限生成 $k$-代数。如果 $k[\xi_1,\ldots,\xi_n]$ 是整环，就重新得到嵌入 $k^n$ 的不可约代数簇的“一般点”。如果 $k[\xi_1,\ldots,\xi_n]$ 是域 $k$ 上的约化代数，仍可给出经典的解释（参见课程提要）；一般情形则完全没有经典解释。现在也就清楚，为什么我们不假设 $I$ 有限：我们要的是全部 $k$-代数构成的范畴，而不只是有限生成代数。$I$ 有限时才得到有限维的标准仿射空间。}，可表函子 $\mathfrak X_A$ 都与某个 $V_S$ 等价。对 $A,B\in\Ob\G_k$，$\mathfrak X_A$ 与 $\mathfrak X_B$ 等价，当且仅当 $A$ 与 $B$ 同构。

\textbf{可表函子 $\mathfrak X_A$（$A\in\Ob\G_k$）是一个范畴 $\Aff_k$ 的对象；这一范畴与反范畴 $\G_k^{\op}$ 等价。}

［$\Aff_k$ 中的态射，按照标准的范畴论方式（见第 1 节），定义为函子态射，即自然变换。它们就是由 $\G_k$ 中态射 $f\colon B\to A$ 诱导的映射 $f^*$：］
\[
\Hom_{k\text{-alg}}(A,k')\longrightarrow\Hom_{k\text{-alg}}(B,k').
\]
这个映射通过与 $f$ 复合来定义：$B\xrightarrow{f}A\xrightarrow{u}k'$，即 $u\mapsto u\circ f$，其中 $u\in\mathfrak X_A(k')$，$k'\in\Ob\G_k$\footnote{进一步的范畴论说明见第二章第 1 节。显然，这里没有使用范畴 $\G_k$ 的任何特殊性质。}。

用 Grothendieck 的原话说：
% PDF 40
\sourcepage{36}{40}
如果 $\mathfrak X$、$\mathfrak X'$ 是分别由 $k$-代数 $A$、$A'$ 表示的 $k$ 上仿射代数空间，那么 $\Aff_k$ 中的态射 $\mathfrak X\to\mathfrak X'$，与方向相反的 $k$-代数同态一一对应：
\[
\Hom_{\Aff_k}(\mathfrak X,\mathfrak X')\xrightarrow{\sim}\Hom_{k\text{-alg}}(A',A).
\]
这只是可表函子的通常推论：函子间的箭头，对应于表示对象间的箭头。（参见上一脚注。）

\textbf{评注。}这些事实看起来过于形式化，甚至像是同义反复。真正有意义的内容，将来自以函子方式把每个 $A$，也即每个 $\mathfrak X_A$，对应到一个几何对象：由 $A$ 定义的仿射概形，即某个局部赋环空间 $(\Spec A,\widetilde A)$（参见第三章）。

目前我们得到的是范畴等价 $\Aff_k\simeq\G_k^{\op}$。当 $k=\mathbb Z$ 时，得到“绝对”仿射代数空间的范畴 $\Aff=\Aff_{\mathbb Z}$，它与含幺交换环的范畴 $\G=\G_{\mathbb Z}$ 反等价。我们认为这是一个几何理论，因为其中所有构造都有几何意义。

如果 $\mathfrak X'\to\mathfrak X$ 是 $\Aff_k$ 中的态射，那么 $\mathfrak X'$ 的每个 $k'$ 值点都被送到 $\mathfrak X$ 的一个 $k'$ 值点（见第 2 节）；这只须对同态作拉回。特别地，$k'=k$ 时也是如此。$f\colon\mathfrak X'\hookrightarrow\mathfrak X$ 是单射，当且仅当对应的 $f^*\colon A\to A'$ 为满射\footnote{译注：原文此处把“单射”与“闭嵌入”混用。满射环同态确实给出闭嵌入，但一般的函子单态并不一定对应满射环同态；例如局部化给出的开嵌入。这里应结合前文由商代数定义的闭代数子空间来理解。}。于是又回到了最初的说明：$f^*$ 表示从 $\mathfrak X$ 到 $\mathfrak X'$ 的限制。特别地，可以像前面那样选择嵌入 $\mathfrak X\hookrightarrow\mathbb E^I$。

在经典情形中，人们特别关心 $k'$ 为约化代数的情况，即 $\operatorname{Nil}k'=0$，或者 $k'$ 是整环，特别是域的情况。此外，如果愿意再次使用我们屡次试图避免的仿射嵌入，那么 $\mathfrak X(k')$ 就是某个 $\mathbb E^I$ 中的闭代数子集。
% PDF 41
\sourcepage{37}{41}
接下来几节将研究 $\mathfrak X$ 的这些特殊类型的“点”，特别是表示 $\mathfrak X$ 的环 $A$ 也属于上述特殊类型的情形。当 $k'$ 为域时，“几何点”（见下面第 9 节）将使我们重新获得经典情形的直观。

\PartTitle{第二部分}{限制到特殊类型的 $k$-代数}
\begin{center}（$k'=k$ 固定、$k'=\mathbb C$、$k'$ 为域、$k'$ 约化，等等）\end{center}
\Section{9}{提要}
\textbf{［}回到多项式方程组 $S$：$f_j(T_i)=0$（$i\in I$，$j\in J$），其系数属于基环 $k$（见第 1 节）。现在不取任意测试代数 $k'\in\Ob\G_k$，而把 $k'$ 限制在 $\G_k$ 的某些子范畴中。这些情形或因几何理由，或因历史原因，或因其对代数几何基础的重要性而值得特别关注；它们也有助于理解 Grothendieck 方法的简化之处。主要情形如下。
\begin{enumerate}[label=\arabic*)]
\item 固定 $k'=k$，也就是固定只有一个对象的子范畴 $\operatorname{Fix}k$，只允许 $k$ 值点。这又回到了算术的观点，参见第 0 节。
\item $k'\in\Ob\mathfrak F_k$，其中 $\mathfrak F_k$ 是 $\G_k$ 中以域为对象的满子范畴。
\item $k$-代数 $k'$ 约化，即 $k'$ 中没有非零幂零元。仍可以像前面一样，把任意 $S$ 化归为 $P_I$ 的理想 $\mathfrak J$；但 $\mathfrak J$ 与 $V_{\mathfrak J}$ 的关系已不再像一般情形那样简单。
\end{enumerate}
在情形 1）中，如果 $k$ 是复数域，那么 $V_S$ 就成为复代数簇 $V_S(\mathbb C)\subset\mathbb C^I$。当 $I$ 有限时，这是十九世纪先驱们——Riemann、Max Noether，以及后来的意大利学派等——的主要研究对象之一\footnote{仿射簇 $V_S(\mathbb C)$ 用于局部问题；对于整体问题，经典代数簇总是指射影簇。}。不可约的仿射或射影代数簇的研究，曾被看作经典代数几何的核心对象。下面沿用 Grothendieck 的原话。\textbf{］}

……现在来看看，如果不取任意代数 $k'$ 和 $k$，而只取某些类型的代数，会发生什么。例如，在最初那个朴素的观点中，我们把自己限制在 $k'=k$ 且 $k'$ 固定的情形（参见情形 1）。起初可以假定 $k'=k=\mathbb R$，即实数域；下一步则是在复数域 $\mathbb C$ 上研究代数几何，并且让它固定不变。
% PDF 42
\sourcepage{38}{42}
这就是经典情形\footnote{参见提要中回顾的经典情形。}，见第 10 节。也可以只取那些本身为域的 $k$-代数 $k'$，或者只考虑约化代数 $k'$\footnote{即没有非零幂零元，参见原书第 9 页脚注 (1)。}。

\Section{10}{域值点}
\textbf{提要。}把函子 $V_{\mathfrak J}$ 限制到 $\G_k$ 的满子范畴 $\mathfrak F_k$，其对象是同时为 $k$-代数的域。准确地说，对 $P_I$ 的每个理想 $\mathfrak J$，定义限制函子
\[
V'_{\mathfrak J}=V_{\mathfrak J}|_{\mathfrak F_k}\colon\mathfrak F_k\longrightarrow\Sets.
\]
此时，映射 $\mathfrak J\mapsto V_{\mathfrak J}$ 的双射性质不再成立；所得结论是：$P_I$ 的两个理想 $\mathfrak J$、$\mathfrak J'$ 定义同一函子，当且仅当二者具有相同的根理想，亦即
\[
V'_{\mathfrak J}\sim V'_{\mathfrak J'}\quad\Longleftrightarrow\quad\sqrt{\mathfrak J}=\sqrt{\mathfrak J'}.
\]
因此，函子 $V'$ 与根理想之间有双射
\[
\mathfrak J=\sqrt{\mathfrak J}\longleftrightarrow V'_{\mathfrak J}\;(=V'_{\sqrt{\mathfrak J}}).
\]
证明需要引理 $\bigcap_{\mathfrak p\in\Spec A}\mathfrak p=\operatorname{Nil}A$，它说明 $A$ 的幂零根的两种定义等价。

$P_I$ 的根理想集合与函子 $V'_{\mathfrak J}$ 之间双射的证明，还能推广到其他情形，例如 $k'$ 约化，特别是 $k'$ 为整环时。不过，由于几何点的重要性，对域的这一结论尤其值得关注，参见下一节。

……所以，如果考察仿射空间 $k'^I$ 中的点，并要求 $k'$ 是域，即 $k'\in\Ob\mathfrak F_k$，其中 $\mathfrak F_k$ 是 $\G_k$ 的满子范畴，那么就是把函子 $V_{\mathfrak J}\colon\G_k\to\Sets$ 限制为
\begin{equation}\tag{10.1}V'_{\mathfrak J}=V_{\mathfrak J}|_{\mathfrak F_k}\colon\mathfrak F_k\longrightarrow\Sets.
\end{equation}
它从同时为 $k$-代数的域所成的范畴 $\mathfrak F_k$ 到集合范畴。这样，我们所考察的就是仿射空间 $\mathbb E^I(k')$ 中坐标属于域 $k'$ 的点。
% PDF 43
\sourcepage{39}{43}
这里 $k'\in\Ob\mathfrak F_k$。若 $f\in P_I$ 且对某个正整数 $m$ 有 $(f(x))^m=0$，其中 $x\in k'^I$，则 $f(x)=0$。这说明，\textbf{如果一个多项式的某次幂在 $V'_{\mathfrak J}(k')$ 的所有点上消失，那么 $f$ 本身也在 $V'_{\mathfrak J}(k')$ 上恒为零}。因此，用根理想 $\sqrt{\mathfrak J}$\footnote{参见原书第 9 页脚注 (2)。}替换 $\mathfrak J$，其中 $\sqrt{\mathfrak J}\supset\mathfrak J$，就有
\begin{equation}\tag{10.2}V'_{\mathfrak J}\sim V'_{\sqrt{\mathfrak J}}.\footnotemark
\end{equation}
\footnotetext{换言之，把 $\G_k$ 限制到 $\mathfrak F_k$ 以后，不仅可以取 $S$ 中多项式的所有线性组合所成的理想 $\mathfrak J$，还可以进一步加入这样的多项式 $f$：它的某个适当的正整数次幂 $f^m$ 已是方程。}

取与 $\mathfrak J$ 关联的理想 $\sqrt{\mathfrak J}$，就是在商环 $P_I/\mathfrak J$ 中取幂零根，再把它拉回到 $P_I$。因此，用这些商环来表述，我们所做的就是把 $P_I/\mathfrak J$ 再对其幂零根 $\operatorname{Nil}(P_I/\mathfrak J)$ 取商，所得商环约化：
\begin{equation}\tag{10.3}P_I/\sqrt{\mathfrak J}\simeq(P_I/\mathfrak J)/\operatorname{Nil}(P_I/\mathfrak J).
\end{equation}
现在看看这意味着什么。容易验证，$V'_{\mathfrak J}\subset V'_{\mathfrak J'}$ 当且仅当 $\sqrt{\mathfrak J'}\subset\sqrt{\mathfrak J}$。接受这一事实，就得到 $P_I$ 的根理想集合与由多项式方程组给出的子函子 $V'_{\mathfrak J}\colon\mathfrak F_k\to\Sets$ 之间反转次序的一一对应。

要建立这一对应，必须说明怎样用 $V'_{\mathfrak J}$ 恢复 $\sqrt{\mathfrak J}$。现在把问题改用交换代数的语言表述。考察商环 $A=P_I/\mathfrak J$。回忆 $V'_{\mathfrak J}$ 的 $k'$ 值点 $u$，就是同态
% PDF 44
\sourcepage{40}{44}
\begin{equation}\tag{10.4}u\colon A\longrightarrow k',\qquad k'\in\Ob\mathfrak F_k.
\end{equation}
于是，刚才叙述的性质可以改述为：$P_I$ 中的元素 $f$ 在每一个形如 (10.4) 的同态下的像都是零。因此，我们所要证明的，就是所有这类同态之核的交恰为 $\sqrt{\mathfrak J}$\footnote{译注：严格说来，$\sqrt{\mathfrak J}$ 是这些核在 $P_I$ 中的原像之交；在 $A$ 中，核的交是 $\operatorname{Nil}A$。}。这个结论等价于交换代数中的一个著名引理\footnote{若环 $A\in\Ob\G$ 的元素 $f$ 满足：对每个到域的同态 $u\colon A\to k$ 都有 $u(f)=0$，则 $f$ 幂零。这是交换代数中的著名结论，和许多同类结果一样，可以形式地从 Krull 定理推出：若 $A\ne0$，则 $A$ 至少含有一个极大理想。}：形如 (10.4) 的同态的核，恰好就是 $A$ 的所有素理想\footnote{$A\in\Ob\G$ 的理想 $\mathfrak p$ 为素理想，当且仅当 $A/\mathfrak p$ 是整环，参见原书第 9 页脚注 (4)。特别地，$A$ 自身绝不是素理想，因为 $A/A\simeq0$ 不是整环；$A$ 的零理想为素理想，当且仅当 $A$ 是整环。}。换言之，这个引理说：
\begin{quote}\bfseries $A$ 的所有素理想之交，等于 $A$ 中幂零元组成的理想，即 $A$ 的幂零根：
\[\bigcap_{\mathfrak p\in\Spec A}\mathfrak p=\operatorname{Nil}A.\]
\end{quote}
如果 $f$ 幂零，那么对每个 $\mathfrak p\in\Spec A$ 都有 $f\in\mathfrak p$，这是显然的。反向蕴含则不那么直接，需要选择公理或 Zorn 引理。

证明如下：若 $f\in A$ 不幂零，则对乘法集 $S=\{f^n\mid n\geq0\}$ 的局部化（参见第三章第 8 节）$S^{-1}A$ 是非零环。因此由 Krull 定理，存在极大理想 $\mathfrak m\subset S^{-1}A$。它在典范同态 $p\colon A\to S^{-1}A$ 下的原像 $\mathfrak p=p^{-1}(\mathfrak m)\in\Spec A$，就是一个不包含 $f$ 的素理想。
% PDF 45
\sourcepage{41}{45}
\Section{11}{经典情形：固定的代数闭域 $k$}
\textbf{提要。}……例如，设 $k'\in\Ob\mathfrak F_k$ 固定且等于 $k$。如果 $k$ 代数闭，而且方程组的不定元集合 $I=\{T_1,T_2,\ldots,T_n\}$ 有限，就能得到一些优美的结果。

将 $k^I$ 中使理想 $\mathfrak J\subset P_I$ 的所有多项式都消失的点集，即 $S$ 的 $k$ 值解集，记为 $V''_{\mathfrak J}$。此时 $V''_{\mathfrak J}$ 仍然只依赖于 $\sqrt{\mathfrak J}$：
\begin{equation}\tag{11.1}V''_{\mathfrak J}=V''_{\sqrt{\mathfrak J}},\qquad\mathfrak J\subset k[T_1,\ldots,T_n],\quad k\text{ 代数闭},\ n<\infty.
\end{equation}
因此，如果 $\mathfrak J$、$\mathfrak J'$ 是 $k[T_1,\ldots,T_n]$ 中具有相同根理想的两个理想，那么相应的 $V''_{\mathfrak J}$、$V''_{\mathfrak J'}$ 相等：
\begin{equation}\tag{11.2}\sqrt{\mathfrak J}=\sqrt{\mathfrak J'}\quad\Longrightarrow\quad V''_{\mathfrak J}=V''_{\mathfrak J'}\subset k^n.
\end{equation}
实际上，知道了 $k^n$ 的闭代数子集 $V''_{\mathfrak J}$，就能重建 $\sqrt{\mathfrak J}$。这实质上就是所谓 Hilbert 零点定理。换言之，凭借 $V''_{\mathfrak J}$ 可以恢复的，并非理想 $\mathfrak J$ 本身，而是它的根理想 $\sqrt{\mathfrak J}$：证明 $\sqrt{\mathfrak J}$ 恰好是在 $V''_{\mathfrak J}$ 上消失的所有多项式 $f\in k[T_1,\ldots,T_n]$ 所成的理想即可。经典上，$V''_{\mathfrak J}$ 称为理想 $\mathfrak J\subset k[T_1,\ldots,T_n]$ 在仿射空间 $k^n$ 中定义的\textbf{仿射代数簇}。

也就是说，我们得到下述经典情形：\textbf{$k[T_1,\ldots,T_n]$ 的任意理想 $\mathfrak J$ 都定义一个仿射代数簇 $V''_{\mathfrak J}(k)=V''_{\mathfrak J}$，其中 $k$ 固定；在 $V''_{\mathfrak J}$ 的每个点上消失的多项式组成的理想，恰好就是 $\mathfrak J$ 的根理想 $\sqrt{\mathfrak J}$。}

这就是 Hilbert 零点定理的经典形式。接下来看看它用内蕴语言如何表述。
% PDF 46
\sourcepage{42}{46}
设 $\mathfrak J$ 是 $P_I=k[T_1,T_2,\ldots,T_n]$ 的理想。仍可将 $k^n$ 的点与 $k$-代数同态 $u\colon P_I\to k$ 认同（见第 2 节）。满足方程组 $\{f(x)=0\mid f\in\mathfrak J\}$ 的那些点，具有 $\ker u\supset\mathfrak J$ 的性质，因此经由商 $A=P_I/\mathfrak J$ 分解，使下图交换：
\[
\begin{tikzcd}[ampersand replacement=\&]P_I\arrow[r,"u"]\arrow[d,"p"']\&k\\A\arrow[ur,dashed]\&\end{tikzcd}
\]
上述结论就是说，如果多项式 $f\in P_I$ 属于所有满足 $\ker u\supset\mathfrak J$ 的代数同态 $u\colon P_I\to k$ 之核的交，那么 $f$ 的某次幂 $f^m$（$m\geq1$）属于 $\mathfrak J$，即\textbf{$f$ 在 $A$ 中的像是幂零元}。由于 $A$ 是有限生成 $k$-代数，而 $k$ 代数闭，零点定理说：
\[
\operatorname{Nil}A=\{f\in A\mid f^m=0\text{，某个 }m\in\mathbb Z_{>0}\}=\bigcap_{u\in\Hom_{k\text{-alg}}(P_I,k)}\ker u.
\]\footnote{译注：原文交式的指标写作 $\Hom_{k\text{-alg}}(P_I,k)$。按这里在环 $A$ 中的表述，应为 $u\in\Hom_{k\text{-alg}}(A,k)$；若用 $P_I$，还需加上 $\ker u\supset\mathfrak J$，并在取交后对 $\mathfrak J$ 取商。}
这一结论可分解为以下两个命题，均针对代数闭域 $k$ 上的有限生成代数 $A$：
\[
\text{零点定理}\quad\Longleftrightarrow\quad
\begin{cases}
\text{a）若 $A$ 是域，则 $A\simeq k$；}\\[3pt]
\text{b）$\displaystyle\operatorname{Nil}A=\bigcap_{\mathfrak m\in\Max A}\mathfrak m$。}
\end{cases}
\]\footnote{$\Max A=\{\mathfrak m\mid\mathfrak m\text{ 是 }A\text{ 的极大理想}\}$，称为 $A$ 的极大谱。}
a）并非同义反复，因为我们的假设并不蕴含 $A$ 是有限生成 $k$-模。零点定理的这一部分说明，要使 a）成立，$k$ 代数闭的假设是必不可少的；否则，$k$ 的任何非平凡有限扩域都会使 a）不成立。

性质 b）说的是，已知为 $A$ 的全部素理想之交的 $\operatorname{Nil}A$，由于我们对 $A$ 的强假设，也等于只对极大理想取交所得的那个原本可能更大的理想。
% PDF 47
\sourcepage{43}{47}
b）是域上有限生成代数的一种十分特殊的性质；一般情况下，它会“错得很厉害”，例如局部环只有一个极大理想。

正向蕴含是清楚的。若 $A$ 为域且 $u\colon A\to k$ 是 $k$-代数同态，那么 $A$ 分解为 $k+\ker u$；但 $\ker u=0$，所以 $A\simeq k$，这给出 a）。b）则来自下述事实：$A$ 的极大理想 $\mathfrak m$ 的特征是 $A/\mathfrak m$ 为域。而 $A/\mathfrak m$ 又是有限生成 $k$-代数，$k$ 代数闭，所以 $A/\mathfrak m\simeq k$。因此，极大谱 $\Max A$ 与集合 $\Hom_{k\text{-alg}}(A,k)$ 之间有双射
\[
\Max A\xrightarrow{\sim}\Hom_{k\text{-alg}}(A,k).
\]
反过来，若 a）、b）成立，那么 $A$ 中所有极大理想之交，就等于从 $A$ 到 $k$ 的全部 $k$-代数同态之核的交，于是得到零点定理。

\textbf{评注。}我们已经指出，a）不能推广到非代数闭域；但 \textbf{b）对任意域 $k$ 都成立}。也就是说，若 $A$ 是域 $k$ 上的有限生成代数，不要求 $k$ 代数闭，那么 $A$ 的极大理想之交仍等于 $A$ 的幂零根。从 $k$ 代数闭的情形不难推出这一结论。a）可替换为 a'）：

\textbf{a'）}若有限生成 $k$-代数 $A$ 是域，那么 $A$ 是 $k$ 的有限代数扩域，即 $A$ 也是有限生成 $k$-模。
% PDF 48
\sourcepage{44}{48}
a'）还可以如下改述：如果 $f$ 属于所有从 $A$ 到某个有限生成扩域的 $k$-代数同态之核的交，那么 \textbf{$f$ 幂零}。

关于零点定理的证明，可参阅 Zariski--Samuel 或 Bourbaki 的《交换代数》。

\Section{12}{$\mathfrak X$ 的点的等价类；几何点}
首先，考虑一个 $k$-代数 $A$，令 $\mathfrak X$ 为由 $A$ 表示的仿射代数空间，即可表函子 $\mathfrak X=V_A$，它给每个 $k$-代数 $k'$ 赋予仿射代数集 $\Hom_{k\text{-alg}}(A,k')$。我们要在 $k'$ 变化时的各个 $k'$ 值点 $P\colon A\to k'$ 之间引入一个等价关系。

一个 $k'$ 值点，也就是 $k$-代数同态
\[P\colon A\longrightarrow k',\]
未必为满射。它给出满射 $P_S\colon A\to P(A)$，其中 $P(A)$ 是 $k'$ 的一个 $k$-子代数，并且可能严格小于 $k'$。对满射 $P_S$，有 $P(A)\simeq A/\ker P_S$，因此可以将 $P_S$ 与典范映射 $a\mapsto a+\ker P_S$ 认同；这一映射仅由核决定。

我们所考虑的等价关系定义为
\[
P\sim P'\quad\Longleftrightarrow\quad\ker P=\ker P'.
\]
等价地，$P\sim P'$ 当且仅当相应的满射点 $P_S$ 与 $P'_S$ 等价。换言之，\textbf{$A$ 的这些点的等价类集合，与 $A$ 的所有理想组成的集合一一对应}；也就是与全部同态像 $A/\mathfrak a$ 连同其典范满同态 $P\colon A\to A/\mathfrak a$ 一一对应。
% PDF 49
\sourcepage{45}{49}
我们特别关心所谓\textbf{几何点}，即\textbf{取值于某个本身为域的 $k$-代数 $k'$ 的点}。这时，$P$ 为几何点，当且仅当 $P(A)\subset k'$ 是整环，亦即 $\ker P_S=\ker P$ 为素理想 $\mathfrak p$\footnote{译注：这里的“当且仅当”在点的等价类意义下理解：$P(A)$ 为整环等价于同核的等价类具有一个域值代表。}。反过来，每个素理想 $\mathfrak p\subset A$ 都通过典范投影 $A\to A/\mathfrak p$ 定义一个几何点的等价类。在由 $\mathfrak p$ 定义的几何点中，有一个“最小”的点，它取值于整环 $A/\mathfrak p$ 的分式域 $K(\mathfrak p)$。$A/\mathfrak p$ 自身为域，当且仅当 $\mathfrak p$ 为 $A$ 的极大理想\footnote{参见原书第 41 页有关素理想的脚注。按照旧时的一种定义，只要求 $a\notin\mathfrak p$、$b\notin\mathfrak p$ 蕴含 $ab\notin\mathfrak p$，那么 $A$ 自身也会被视为素理想。这里必须记住，我们遵循 Grothendieck 的约定，因为集合 $\Spec A=\{\mathfrak p\mid\mathfrak p\text{ 为 }A\text{ 的素理想}\}$，即 $A$ 的谱，有重要作用（见第三章）。因此 $\Spec A=\varnothing\Longleftrightarrow A=0$，这是 Krull 定理的推论。}。

\Section{13}{关于幂零元的评议}
\textbf{Grothendieck 的讲述。}

……所以我们看到，在某些问题中，没有把所有 $k$-代数 $k'$ 都作为“测试代数”，即作为代数空间之点的坐标所取值的代数，而只选取其中某些类型：域、整环、约化环，或者 $k$ 的有限扩域；若 $k$ 代数闭，则固定 $k'=k$。这样做的代价，是无法再从几何上区分 $P_I$ 的理想 $\mathfrak J$ 与其根理想 $\sqrt{\mathfrak J}$。换言之，［再次摆脱仿射嵌入，］从内蕴的函子观点出发，考察由 $A=P_I/\mathfrak J$ 表示的函子 $\mathfrak X_A$，那么由于对 $k'$ 作了限制，就无法区分 $A$ 与其商 $A/\operatorname{Nil}A$；后者满足 $\operatorname{Nil}(A/\operatorname{Nil}A)=0$。所以，只用域或某些类型的域 $k'$ 进行研究，在技术上就意味着只用约化环 $A$［作为 $\mathfrak X_A$ 的表示对象］。

这种做法持续了很长时间，直到十年或十五年以前仍是如此！当然，研究交换代数的人一直会考察非约化环，［例如 $\mathbb Z/p^n\mathbb Z$、$P_I/\mathfrak p^n$ 等，］但在几何中，人们在某种意义上拒绝给非约化环赋予几何解释。真是一种令人不快的处境！Kähler 第一个系统地建立了这样的框架：把几何对象也赋予可能含幂零元的环。Dieudonné 和我沿着同一方向继续前进。

过去十五年的大量进展已清楚表明，在代数几何中，只研究约化环，无法得到我们所期待的全部结果。在许多问题中，\textbf{幂零元恰恰是关键所在}。例如在无穷小问题中，幂零元就起着极为重要的作用。

\begin{quote}\bfseries 因此，我认为应当坚持如下的一般原则：不对 $k$ 上的代数 $k'$ 作任何限制，当然也不假定它们约化；也就是说，绝不为了消去幂零元而把 $k$-代数 $A$ 与 $A/\operatorname{Nil}A$ 混为一谈，因为幂零元携带着关于环 $A$ 的极有价值的信息。\end{quote}
% PDF 50
\sourcepage{46}{50}
\textbf{［}我想解释一下，为什么经典代数几何学家不考虑幂零元。若 $k$ 代数闭，经典上取 $k=\mathbb C$，当时所考虑的仿射代数空间，也即仿射簇，被认同为极大谱 $\Max A$。于是，对每个 $f\in A$，可以赋予函数
\[
\widetilde f\colon\mathfrak m\longmapsto\widetilde f(\mathfrak m)
=\text{$f$ 在 $A_{\mathfrak m}$ 的剩余域（$\simeq k$）中的像}.
\]
这一函数可以看作“多项式函数”。映射 $f\mapsto\widetilde f$ 是环同构，因为 $\widetilde f$ 恒为零，当且仅当对每个 $\mathfrak m\in\Max A$ 都有 $f\in\mathfrak m$；这蕴含 $f$ 幂零，进而 $f=0$，因为 $A$ 约化（参见第 11 节，原书第 42 页）。

\textbf{经典上，$A$ 被看作 $\mathfrak X_A(k)$ 上的多项式函数环}；这些函数在 Zariski 拓扑下连续，参见下一节。\textbf{如果 $A$ 不约化，这种函数解释就失效了。}如果坚持把非零幂零元 $f\in A$ 称为“函数”，就会遇到看似矛盾的情形：$f$ 非零，而 $\widetilde f$ 却恒为零。如果没有层的发明，就难以克服这一反对意见。

Serre 在 FAC 中证明，$A$ 是结构层 $\mathcal O_X$ 的整体截面环，其中 $X=\Max A$，$A$ 约化；$\mathcal O_X$ 在每个 $\mathfrak m\in\Max A$ 处的茎就是局部环 $A_{\mathfrak m}$。去掉对 $\operatorname{Nil}A$ 的限制以后，仍能构造 $\mathcal O_X$，并仍有 $A=\Gamma(X;\mathcal O_X)$；但是，把 $\mathcal O_X$ 的元素称为函数芽就不正确了。

第三章中，Grothendieck 将说明，对于任意含幺交换环 $A\in\Ob\G$（见第 2 节），如何用全谱 $\Spec A$ 替换 $\Max A$，从而完成全部构造：Zariski 拓扑以及层 $\mathcal O_X$。
% PDF 51
\sourcepage{47}{51}
\textbf{只用极大谱无法得到令人满意的结果……。］}

\Section{14}{Zariski 拓扑}
\textbf{提要。}设 $k$ 为代数闭域。$k^n$ 中所有代数簇组成的族 $\mathcal V(k^n)$ 包含空集 $\varnothing$ 和 $k^n$，并且对有限并和任意交封闭。因此，$\mathcal V(k^n)$ 是 $k^n$ 上某个拓扑的闭集族；这一拓扑称为 Zariski 拓扑。

假定方程组 $S$ 只由一个常数多项式 $c$ 组成。$c=0$ 时，$k^n$ 的每个点都是解；$c\ne0$ 时，$k^n$ 没有点是解。因此 $\varnothing,k^n\in\mathcal V(k^n)$（参见提要）。

设 $\{V_\nu\}_{\nu\in N}$ 是 $k^n$ 中的一族代数簇，分别由多项式方程组 $S_\nu$（$\nu\in N$）表示。令 $S=\bigcup_{\nu\in N}S_\nu$ 为这些方程的并，那么 $S$ 所表示的就是 $V_\nu$ 的集合论交：
\begin{equation}\tag{14.1}V_S=\bigcap_{\nu\in N}V_\nu.
\end{equation}
上述说明事实上对任意含幺交换环 $k$ 都成立。现在若 $k$ 为整环，$S_1$、$S_2$ 为两个多项式方程组，那么立即可见，方程组
\begin{equation}\tag{14.2}S=\{f_1f_2\mid f_1\in S_1,\ f_2\in S_2\}
\end{equation}
表示集合论并：
\begin{equation}\tag{14.3}V_S=V_{S_1}\cup V_{S_2}.
\end{equation}
事实上，$V_{S_1}\cup V_{S_2}\subset V_S$ 显然对任意 $k$ 都成立。另一方面，若 $x\in V_S$，则
% PDF 52
\sourcepage{48}{52}
\begin{equation}\tag{14.4}(f_1f_2)(x)=f_1(x)f_2(x)=0\quad\Longrightarrow\quad f_1(x)=0\ \text{或}\ f_2(x)=0,
\end{equation}
因为 $k$ 是整环。利用归纳法，性质 (14.3) 可以推广到任意有限并。

这些简单性质在经典时期早已为人熟知\footnote{例如参见 Hodge--Pedoe，\textit{Methods of Algebraic Geometry}［《代数几何方法》］，Cambridge University Press，1947，第 II 卷。书中验证了正文列出的性质，却没有指出这些性质刻画了一个拓扑的闭集。顺便说，同样的性质也适用于射影空间，只须将 $S$ 中的多项式改为齐次多项式。}。然而，多年来 Zariski 拓扑没有受到认真重视，因为它具有下面这些当时看来十分不理想的性质：
\begin{enumerate}[label=\arabic*)]
\item 它非常粗！
\item 闭集非常“薄”。等价地，开集“非常大”（2'）。
\item 任意非空开集的闭包都是整个空间。
\item 任意两个非空开集 $U$、$V$ 都相交：$U\cap V\ne\varnothing$！
\end{enumerate}
性质 4）意味着，只要 $n>0$，Zariski 拓扑就不是 Hausdorff 的。不过，它满足 $T_1$ 分离公理，即每个单点集都是闭集，因为它是 $n$ 个线性无关的超平面的交！

Zariski 拓扑比 $\mathbb C^n$ 的通常拓扑更粗。早期研究者以这种自然拓扑为基础来运用拓扑方法，例如将其与解析函数理论中的超越方法相结合。

Serre 的 FAC 论文（1955）用 Zariski 拓扑替代解析流形的自然拓扑，表明此前用于解析凝聚层的许多方法，也可以推广到纯代数的情形。
% PDF 53
\sourcepage{49}{53}
第三章将说明，在更一般的情形中，这种最初朴素意义下的 Zariski 拓扑仍然有意义：点变成某个含幺交换环 $A$ 的素理想。实际上，(14.3) 自然地变为
\begin{equation}\tag{14.3$'$}
\text{“$(f_1f_2)(\mathfrak p)=0$”}\quad\Longrightarrow\quad
\text{“$f_1(\mathfrak p)=0$”或“$f_2(\mathfrak p)=0$”}.
\end{equation}
这里约定用 $f(\mathfrak p)=0$ 表示 $f\in\mathfrak p$，或者说 $f\equiv0\pmod{\mathfrak p}$。

\BookChapter{2}{仿射代数空间范畴中的极限}
\begin{quote}
［初读时可以略去第二章。Grothendieck 讲授代数群课程时需要这一章的内容。对任意含幺交换环 $A$ 构造其仿射概形 $(\Spec A,\widetilde A)$，不需要第二章的任何内容。建议初学者先读第三章。］
\end{quote}

……我们已建立了关于仿射代数空间的基本语言［参见第一章第 8 节］，接下来可以沿两个方向前进：其一，以另一种方式重新表述这个概念，使得我们能够把仿射片粘合起来，定义不必仿射的更一般代数空间，也就是不必由 $k$-代数表示的空间\footnote{……这引导我们定义谱 $X=\Spec A$ 及结构层 $\mathcal O_X$，然后把若干仿射片粘合起来［第二卷］。}；其二，先在这些受限对象组成的范畴 $\Aff_k$ 内多做一些工作，定义积、纤维积、核等运算［参见第 1、6 节］。由于接下来将引入仿射群对象\footnote{这些构造有助于在范畴 $\Aff_k$ 中定义群对象；这正是他同时在 SUNY 布法罗分校讲授的仿射代数群课程所需的内容，1973 年夏，参见 AAG--Buffalo course。}，即 $\Aff_k$ 中的群对象，我们先选第二条路线，然后再将这些概念整体化［见第二卷］。

\Section{0}{提要}
\textbf{［}第二章主要研究 $k$ 上仿射代数空间范畴 $\Aff_k$ 中的极限（参见第 2 节），其中 $k\in\Ob\G$；特别地，$k=\mathbb Z$ 时，就是“绝对情形”$\Aff=\Aff_{\mathbb Z}$。

第一章已经将 $\Aff_k$ 的对象定义为解函子。
